\clearpage\nobalance
\section*{Author Notes}
\subsection*{Acknowledgements}
I thank Prof. Luyao Zhang, Program Chair, and Prof. Ken Rogerson, Program Discussant, of the \emph{2056 Nobel and Turing Laureate Program Launch Workshop}, Duke Kunshan University, September 7, 2026; Longfei Jing and Yongzhuo Wu, my teammates in Workshop 6; and the rest of the COMSCI/ECON 206 Autumn 2026 Session 1 class. I am responsible for the proposal and remaining errors.
\subsection*{Individual contribution}
I formed the core idea, posed the research question, and developed the one-click verifier, the lemons-market framing and the rule triple $(m,p,F)$. I verified the literature, played the game, and on 9 September found and corrected two errors in the payoff law (Appendix~\ref{app:develop}). The photographs, the numbers and the final wording are mine.

\bibliographystyle{ACM-Reference-Format}\bibliography{references}
\appendix
\section{Technical details and reproducibility}\label{app:tech}
\textbf{Notation.} $a_i$ developer $i$'s action (honest, none, fake); $v_i$ verifier enrolment; $m$ mandatory attachment; $p$ spot-check rate; $F$ fine on a detected fake; $pF$ expected fine; $\Delta$ premium paid only for a verifiable trace file; $p^{*}$ smallest $p$ making honesty the unique best response.

\textbf{Artifact.} Trace-or-Fake Market Game, a $3\times2$ simulator. Space: \url{https://huggingface.co/spaces/Peterhe123/Trust}. Notebook: \texttt{colab/trace\_or\_fake\_psweep.ipynb}. Author: Zhengjun He. v1.0.0, 6 September 2026.

\textbf{Payoff law.} A developer chooses honest, none or fake, and whether to enrol in the verifier. Payoff is market benefit plus $\Delta$ if honest, minus cost, plus stability if enrolled, minus $pF$ if faking. A fake earns the same base benefit as an honest deliverable, since buyers cannot observe truth before purchase, while saving cost.

\textbf{Corrections, 9 September 2026.} (i) The first law summed a cost, a benefit and a stability index. Cost and benefit have opposite signs, so summing them scored honesty as better on both before any rule applied, and the sweep returned $p^{*}=0$: faking was never attractive, so the question could not be posed. Cost now enters negatively. (ii) Honest and faked traces were then paid identically, so honesty won only by avoiding a penalty. A premium $\Delta$ paid only for a verifiable trace file repairs this; the two mechanisms substitute, giving $\Delta+pF>25$ against a fake advantage of $25$. At $\Delta=0$, $p^{*}=0.5$ at $F=50$; at $\Delta=15$, $p^{*}=0.2$; at $\Delta=25$, no enforcement is needed.

\textbf{Fresh-run record.} Python 3.11, standard library only, deterministic, no seed needed. Six checks pass: no-rules fake wins; enforcement alone makes honesty unique; a premium alone ($\Delta=30$) also does; status quo unchanged; enrolment weakly dominant; six cells. Abridged output:

\begin{quote}\small\ttfamily
no rules: A 135, B 120, C 160  ->  best response C\_fake\\
threshold p* = 0.5000 (F = 50);  F=100 -> 0.2500;  F=400 -> 0.0625\\
premium sweep (p=F=0): 0 -> C\_fake;  25 -> tie;  30 -> A\_honest\\
combined: delta=0 -> p*=0.50;  delta=15 -> p*=0.20;  delta=25 -> no enforcement
\end{quote}

\textbf{Boundary.} Static client-side HTML/CSS/JavaScript only: no runtime, backend, key or paid service. Hosting and the viewer's browser still consume energy~\cite{lannelongue2021}. Semantic fieldsets, table captions and keyboard operability are present; a full screen-reader pass is unverified. The deployed Space should be updated to the corrected law.

\subsection{AI-use disclosure}
I completed all substantive work myself: the research question, the payoff law, the strategy matrix, the two corrections to that law, the derivation of the threshold $\Delta+pF>25$, the computational checks, the game, the notebook and the argument of the paper. Every source was read and every number re-derived by me, and I played the game and ran the checker. A generative AI assistant (DeepSeek, model \texttt{deepseek-v4-flash}, accessed 6 September 2026) was used only for \LaTeX{} typesetting and layout, and for brainstorming possible literature to look up. It did not devise the model, run any computation, or write the argument. All literature it suggested was independently located and verified before citation. Initial reasoning and handwritten reflection are Human-Only, and I remain responsible for every claim.

\section{Cumulative Intellectual Development}\label{app:develop}
This proposal develops my Intellectual Statement II through the \emph{2056 Nobel and Turing Laureate Program Launch Workshop}, September 7, 2026, Duke Kunshan University, IB 2050. Program Chair: Prof. Luyao Zhang. Program Discussant: Prof. Ken Rogerson. My session: Workshop 6. My role: Computer scientist. My teammates: Longfei Jing and Yongzhuo Wu.

\textbf{First change: the object of the question.} My earlier claim was that I would build mechanisms letting agents cooperate under near-total unobservability. The September 4 evidence forced something narrower: at the Tencent Shanghai Office plugin quality is hard to observe before adoption, and at the Shanghai Science and Technology Museum the deep-residual-learning panel asks a general audience to trust institutional authority rather than checkable evidence. The marketplace is a market without institutional trust; the museum is institutional trust without a market. Workshop feedback sharpened the boundary---build-layer provenance already exists, so the contribution must sit at the claim layer---moving the proposal from a disclosure design to a rule design.

\textbf{Second and third changes, 9 September 2026.} The notebook sweep returned $p^{*}=0$, showing the game summed cost with the wrong sign, so honesty won regardless of any rule and the game could not pose its own question. With that fixed, honest and faked traces were still paid identically, so honesty was a tax rather than an asset. Reading M{\"a}kimattila et al.~\cite{makimattila2025}, where a sender's \emph{choice} among offered tests is itself informative, together with Spence's single-crossing condition~\cite{spence1973}, prompted adding the premium $\Delta$. The proposal can now show honesty wins because it is worth more, not merely because cheating is punished. What remains unresolved: $\Delta$ is still a parameter rather than a belief, so making it endogenous is the next study.

\section{Field and open-source inspiration}
The \emph{Joint Interdisciplinary Field Trip---Technology, Computation, Culture \& Society in Action} took place on September 4, 2026.

\textbf{Tencent Shanghai Office.} Observation: WorkBuddy is Tencent's AI agent platform, whose marketplace connects plugin developers with users. Interpretation: plugin quality is hard to observe before adoption, so buyers rely on cheap signals, and costly verification can produce a lemons market. Photograph: \emph{WorkBuddy agent-plugin marketplace}, by the author, 4 September 2026.

\textbf{Shanghai Science and Technology Museum.} Observation: the 2023 Mathematics and Computer Science Prize panel~\cite{futureprize2023} presents deep residual learning through simplified diagrams. Interpretation: trust here rests on institutional authority, not on independently checkable evidence. Photograph: \emph{2023 Mathematics and Computer Science Prize panel}, by the author, 4 September 2026.

\textbf{Two open-source tools.} (i) \emph{SLSA}~\cite{slsa}, \url{https://slsa.dev/}, v1.2, Apache-2.0: generates non-forgeable build provenance, so a verifier can check which commit built which artifact, but certifies the build chain only, not whether the artifact does what its description claims. (ii) \emph{CKAN}~\cite{ckan}, \url{https://ckan.org/}, 2.11, AGPL-3.0: an open-data catalog with standardised metadata and public APIs that stores whatever publishers upload and never checks whether a dataset matches its description. Consequence: SLSA supplies the build layer at near-zero cost and CKAN the public substrate, so my contribution is the claim layer both ignore.

\section{Review response and revision record}
No peer review has been received as of v1. In-class review is September 14, 2026; no reviewer feedback is claimed. This section will be completed with the author response and the v2 revision.

\clearpage\onecolumn
\section{Structured abstract and cumulative proposal map}\label{app:abstract}
\par\medskip
\begin{table}[H]
\caption{Appendix E structured abstract; the eight-row proposal structure.}
\label{tab:abstract-map}
\renewcommand{\arraystretch}{1.18}
\begin{tabularx}{\textwidth}{@{}p{0.27\textwidth}Y@{}}
\toprule
\textbf{Proposal element} & \textbf{Current statement} \\
\midrule
Background & Build-layer traceability already exists and generation is near-free, while unchecked quality uncertainty collapses markets. Foundations: Akerlof (1970); Buneman et al. (2001). \\
Motivation and application & Untested: which mechanisms make honest trace files an equilibrium when they are free to generate and to fake. Scenario: a plugin marketplace setting the rules, with users and journalists as verifiers; it scales to public claims. \\
Three connected questions & Economics: does a fine $pF$ or a premium $\Delta$ make honesty a Nash equilibrium? Computation: what format and verifier make a per-claim declaration checkable at near-zero cost? Behavior: does traceability shrink the justification space, or does easy checking license less scrutiny? Integrated: which rules make honest disclosure stable when the evidence can be forged? \\
Method and model assumptions & Developers choose honest, none or fake, plus enrolment; the platform commits to $(m,p,F)$ and the market sets $\Delta$. Information is incomplete. Payoff is benefit plus $\Delta$ if honest, minus cost, plus stability if enrolled, minus $pF$ if faking; a fake earns the same base benefit as an honest deliverable. Baseline nets $120$; the computation is a $3\times2$ enumeration plus sweeps of $p$ and $\Delta$. \\
Results & Verified: with no rules and $\Delta=0$ the fake wins ($160$ vs $135$); a premium alone restores honesty, and without one honesty requires $pF>25$ ($p^{*}=0.5000$ at $F=50$). The two substitute: $\Delta=15$ gives $p^{*}=0.20$; $\Delta=25$ needs no enforcement. Six checks pass. Evidence type: simulated, from one shared payoff law. Two corrections; see Appendix~\ref{app:develop}. \\
Intellectual merits & Checkability becomes an institutional variable rather than a technical property, extending lemons theory to claims under free generation, where the certification instrument is itself forgeable. Separating enforcement from reward and showing they are substitutes reframes the design problem from policing to pricing. \\
Practical impacts & Private: lower screening costs on plugin and agent platforms. Public: auditable claims in museums, ministries and data portals. The SDG 4 tool lets students and the public learn why honest disclosure needs cheap verification, assessable by a short pre/post prediction question. The simulator is released; learning gains remain unevaluated. \\
Feedback received and revision made & September 7 workshop feedback established that build-layer provenance already exists, moving the contribution to the claim layer and revising the 2056 statement from cooperation under unobservability to a governed trust infrastructure. September 9, two dated changes: the sweep returned $p^{*}=0$, exposing a sign error; and reading M\"akimattila et al.~\cite{makimattila2025} showed honesty was modelled as a cost rather than an asset, so a premium $\Delta$ was added. Peer reviews pending; none claimed. \\
\bottomrule
\end{tabularx}
\end{table}
